\begin{figure}[H]
\centering
\begin{tikzpicture}[
    node distance=0.9cm and 1.4cm,
    box/.style={draw, rounded corners, minimum width=2.2cm,
                minimum height=0.75cm, align=center, font=\scriptsize},
    arrow/.style={-Latex, line width=0.3pt},
    dashedbox/.style={draw, rounded corners, densely dashed, inner sep=4pt}
]

% Neural model
\node[box, fill=blue!10] (model) {Neural model\\(proposal)};
\node[box, below=0.9cm of model] (sample) {Candidate\\solution};

\draw[arrow] (model) -- (sample);

% CO solver
\node[box, fill=green!10, right=2.6cm of sample] (solver) {CO solver\\(black-box)};
\draw[arrow] (sample.east) -- (solver.west);

% Repair / feedback
\node[box, fill=yellow!20, above=0.85cm of solver] (repair) {Repair / validate};
\draw[arrow] (solver.north) -- (repair.south);

% Feedback loop
\node[box, fill=blue!10, above=0.85cm of model] (update) {Learning update};
\draw[arrow] (repair.west) -- (update.east);
\draw[arrow] (update.south) -- (model.north);

% Group labels
\node[dashedbox, fit=(model)(update), label=left:{\tiny Learning system}] {};
\node[dashedbox, fit=(solver)(repair), label=right:{\tiny CO solver (oracle)}] {};

\end{tikzpicture}

\caption{Illustration of black-box learning with a CO solver. A neural model proposes candidate solutions, which are evaluated and possibly repaired by a conventional optimization solver. The solver's feedback is used to improve the proposal mechanism without requiring differentiability or solver-side modifications.}
\label{fig:blackbox_learning}
\end{figure}